\chapter{Resultados y discusión}\label{cap:resultados}

\section{Comparativa global de modelos}

% TODO(resultados): completar con métricas reales tras Fase 6.
\begin{table}[h]
  \centering
  \begin{tabular}{lcccc}
    \toprule
    Modelo & Acc. & F1 macro & NLL & ECE \\
    \midrule
    TF-IDF + LR             & --- & --- & --- & --- \\
    DeBERTa-v3 base         & --- & --- & --- & --- \\
    DeBERTa-v3 + temp.\ scaling & --- & --- & --- & --- \\
    \bottomrule
  \end{tabular}
  \caption{Métricas finales sobre la partición \emph{dev} de FEVER.}
  \label{tab:resultados-globales}
\end{table}

\section{Análisis por clase}

% TODO(resultados): tabla F1 por clase.
\begin{table}[h]
  \centering
  \begin{tabular}{lccc}
    \toprule
    Modelo & F1 \textsc{sup} & F1 \textsc{ref} & F1 \textsc{nei} \\
    \midrule
    TF-IDF + LR     & --- & --- & --- \\
    DeBERTa-v3      & --- & --- & --- \\
    \bottomrule
  \end{tabular}
  \caption{Desglose por clase del F1.}
  \label{tab:resultados-por-clase}
\end{table}

La clase \textsc{nei} suele resultar la más difícil porque exige
distinguir ausencia de información de presencia de información
contraria; se discutirá si esta dificultad se traduce en menor
calibración para esa clase.

\section{Calibración}

Se incluyen los diagramas de fiabilidad antes y después de
\emph{temperature scaling} para el modelo neuronal. Se espera reducción
visible de ECE sin pérdida de exactitud.

\section{Contrastes estadísticos}

Conforme al capítulo~\ref{cap:tests}:
\begin{itemize}[leftmargin=*,itemsep=0.2em]
  \item Test de McNemar DeBERTa vs.\ baseline:
        $\chi^2 = ---$, $p = ---$.
  \item Bootstrap $\Delta F1_{\text{macro}}$:
        IC 95\,\% $= [\,---,\;---\,]$.
\end{itemize}

\section{Curva precisión--cobertura por umbral $\tau_{\text{conf}}$}

Se reporta, para distintos valores de $\tau_{\text{conf}} \in
\{0.5, 0.55, 0.6, 0.7, 0.8\}$, la fracción de \emph{claims} con
veredicto $\neq$ \textsc{nei} (cobertura) y la precisión sobre los
veredictos emitidos. La elección final de $\tau_{\text{conf}}$
corresponde al punto de operación que maximiza F1 sobre claims con
veredicto.

\section{Análisis de errores}

Se analizan cualitativamente 30 ejemplos donde DeBERTa falla y el
baseline acierta, y viceversa, identificando patrones recurrentes
(negaciones, números, entidades nombradas, comparativos).

\section{Limitaciones}

\begin{itemize}[leftmargin=*,itemsep=0.2em]
  \item FEVER se construye sobre Wikipedia y por tanto el dominio es
        enciclopédico; la generalización a dominio periodístico no se
        evalúa aquí.
  \item El supuesto de independencia condicional de evidencias
        (capítulo~\ref{cap:agregacion}) es una aproximación.
  \item No se evalúa la subtarea de recuperación; se asumen evidencias
        ya disponibles.
\end{itemize}
